% GENERATED FILE. Edit canonical lessons, then run npm run generate:books.
% canonical-source-hash: fnv1a64:160bcf4d22f901db
% canonical-lessons: GE-C275-angeblich, GE-C275-offenbar, GE-C275-zufolge, GE-C275-vermuten, GE-C275-vermutlich

\chapter{Angeblich}
\label{ch:ge-angeblich}
\hlchaptermodality{275}

\begin{chapteropening}
\textbf{By the end of this chapter:} \emph{I can report what others say without promising it is true: supposedly, apparently, according to, to suspect and presumably.}

Complete the last lesson of chapter 275: I can report what others say without promising it is true: supposedly, apparently, according to, to suspect and presumably.
\end{chapteropening}

\section[angeblich]{angeblich — supposedly, allegedly}

\label{lesson:GE-C275-angeblich}



\begin{quote}
Before the new one: say the German for the claim, then the German for to confirm.
\end{quote}

\subsection*{You'll want to know: angeblich}
\textbf{angeblich} — \textquotedblleft{}supposedly, allegedly\textquotedblright{}.

\begin{quote}
Angeblich war der Dieb ein junger Mann.
\end{quote}

\begin{quote}
\emph{Supposedly the thief was a young man.}
\end{quote}

The speaker does not promise that it is true.

\begin{cousinweb}
\emph{angeblich} is built on \emph{angeben} \textquotedblleft{}to state\textquotedblright{}, from \emph{geben} \textquotedblleft{}to give\textquotedblright{}, Old High German \emph{geban} (EtymDB): as it is stated.
\end{cousinweb}

\begin{grammarlens}[title={What you've built}]
One more word for telling, explaining and describing at B1.
\end{grammarlens}

\subsection*{Your turn}
\begin{itemize}
\raggedright
  \item \emph{Say it:} \emph{angeblich}
  \item \emph{Say it:} \emph{angeblich}, once more
  \item \emph{Say it:} a sentence of your own with \emph{angeblich}
  \item \emph{Recall:} say the German for the claim, then the German for to confirm, then say \emph{angeblich} again
\end{itemize}

\begin{tcolorbox}[breakable,colback=teal!4,colframe=teal!35!black,title={Before you move on}]
What does angeblich mean? (\textquotedblleft{}Supposedly, allegedly\textquotedblright{}.) Say it once more.
\end{tcolorbox}
\section[offenbar]{offenbar — apparently, evidently}

\label{lesson:GE-C275-offenbar}



\begin{quote}
Before the new one: say the German for to confirm, then the German for supposedly.
\end{quote}

\subsection*{You'll want to know: offenbar}
\textbf{offenbar} — \textquotedblleft{}apparently, evidently\textquotedblright{}.

\begin{quote}
Er war offenbar sehr müde.
\end{quote}

\begin{quote}
\emph{He was evidently very tired.}
\end{quote}

It sounds sure, because there is a sign: \emph{Offenbar hat es geregnet, die Straße ist nass}.

\begin{cousinweb}
\emph{offenbar} joins \emph{offen} \textquotedblleft{}open\textquotedblright{}, an inherited Germanic word, and the ending \emph{-bar}: plain to see. EtymDB lists no parent for the adverb.
\end{cousinweb}

\begin{grammarlens}[title={What you've built}]
One more word for telling, explaining and describing at B1.
\end{grammarlens}

\subsection*{Your turn}
\begin{itemize}
\raggedright
  \item \emph{Say it:} \emph{offenbar}
  \item \emph{Say it:} \emph{offenbar}, once more
  \item \emph{Say it:} a sentence of your own with \emph{offenbar}
  \item \emph{Recall:} say the German for to confirm, then the German for supposedly, then say \emph{offenbar} again
\end{itemize}

\begin{tcolorbox}[breakable,colback=teal!4,colframe=teal!35!black,title={Before you move on}]
What does offenbar mean? (\textquotedblleft{}Apparently, evidently\textquotedblright{}.) Say it once more.
\end{tcolorbox}
\section[zufolge]{zufolge — according to}

\label{lesson:GE-C275-zufolge}



\begin{quote}
Before the new one: say the German for supposedly, then the German for apparently.
\end{quote}

\subsection*{You'll want to know: zufolge}
\textbf{zufolge} — \textquotedblleft{}according to\textquotedblright{}.

\begin{quote}
Dem Bericht zufolge gab es keine Verletzten.
\end{quote}

\begin{quote}
\emph{According to the report there were no injured people.}
\end{quote}

It follows a dative noun: \emph{dem Bericht zufolge}. Formal and common in the news.

\begin{cousinweb}
\emph{zufolge} joins \emph{zu} and \emph{Folge} \textquotedblleft{}what follows\textquotedblright{}, from Middle High German \emph{volge} (EtymDB): following the report.
\end{cousinweb}

\begin{grammarlens}[title={What you've built}]
One more word for telling, explaining and describing at B1.
\end{grammarlens}

\subsection*{Your turn}
\begin{itemize}
\raggedright
  \item \emph{Say it:} \emph{zufolge}
  \item \emph{Say it:} \emph{zufolge}, once more
  \item \emph{Say it:} a sentence of your own with \emph{zufolge}
  \item \emph{Recall:} say the German for supposedly, then the German for apparently, then say \emph{zufolge} again
\end{itemize}

\begin{tcolorbox}[breakable,colback=teal!4,colframe=teal!35!black,title={Before you move on}]
What does zufolge mean? (\textquotedblleft{}According to\textquotedblright{}.) Say it once more.
\end{tcolorbox}
\section[vermuten]{vermuten — to suspect, to presume}

\label{lesson:GE-C275-vermuten}



\begin{quote}
Before the new one: say the German for apparently, then the German for according to.
\end{quote}

\subsection*{You'll want to know: vermuten}
\textbf{vermuten} — \textquotedblleft{}to suspect, to presume\textquotedblright{}.

\begin{quote}
Ich vermutete, dass er krank war.
\end{quote}

\begin{quote}
\emph{I suspected that he was ill.}
\end{quote}

With \emph{dass} and the verb at the end: \emph{Ich vermute, dass er recht hat}.

\begin{cousinweb}
\emph{vermuten} is built on \emph{Mut}, once \textquotedblleft{}mind, mood\textquotedblright{}, from Proto-Germanic \emph{mōdaz} (EtymDB), with the prefix \emph{ver-}: to have something in mind.
\end{cousinweb}

\begin{grammarlens}[title={What you've built}]
One more word for telling, explaining and describing at B1.
\end{grammarlens}

\subsection*{Your turn}
\begin{itemize}
\raggedright
  \item \emph{Say it:} \emph{vermuten}
  \item \emph{Say it:} \emph{vermuten}, once more
  \item \emph{Say it:} a sentence of your own with \emph{vermuten}
  \item \emph{Recall:} say the German for apparently, then the German for according to, then say \emph{vermuten} again
\end{itemize}

\begin{tcolorbox}[breakable,colback=teal!4,colframe=teal!35!black,title={Before you move on}]
What does vermuten mean? (\textquotedblleft{}To suspect, to presume\textquotedblright{}.) Say it once more.
\end{tcolorbox}
\section[vermutlich]{vermutlich — presumably}

\label{lesson:GE-C275-vermutlich}



\begin{quote}
Before the new one: say the German for according to, then the German for to suspect.
\end{quote}

\subsection*{You'll want to know: vermutlich}
\textbf{vermutlich} — \textquotedblleft{}presumably\textquotedblright{}.

\begin{quote}
Der Zug kommt vermutlich zu spät.
\end{quote}

\begin{quote}
\emph{The train will presumably be late.}
\end{quote}

Close to \emph{wahrscheinlich}. Built on \emph{vermuten}.

\begin{cousinweb}
\emph{vermutlich} is built on \emph{vermuten} \textquotedblleft{}to suspect\textquotedblright{}, which rests on \emph{Mut} \textquotedblleft{}mind\textquotedblright{}, from Proto-Germanic \emph{mōdaz} (EtymDB).
\end{cousinweb}

\begin{grammarlens}[title={What you've built}]
One more word for telling, explaining and describing at B1.
\end{grammarlens}

\subsection*{Your turn}
\begin{itemize}
\raggedright
  \item \emph{Say it:} \emph{vermutlich}
  \item \emph{Say it:} \emph{vermutlich}, once more
  \item \emph{Say it:} a sentence of your own with \emph{vermutlich}
  \item \emph{Recall:} say the German for according to, then the German for to suspect, then say \emph{vermutlich} again
\end{itemize}

\begin{tcolorbox}[breakable,colback=teal!4,colframe=teal!35!black,title={Before you move on}]
What does vermutlich mean? (\textquotedblleft{}Presumably\textquotedblright{}.) Say it once more.
\end{tcolorbox}
